% !TeX encoding = UTF-8
\documentclass[variant=guia,cover=institutional,mode=screen]{ua-engineering-v6-article}
% !TeX encoding = UTF-8
\UASetUniversity{Universidad Austral}
\UASetFaculty{Facultad de Ingeniería}
\UASetDepartment{Departamento de Computación}
\UASetCampus{Pilar}
\UASetCareer{Ingeniería Informática}
\UASetCommission{Comisión A}
\UASetProfessor{Equipo docente}
\UASetSemester{Segundo cuatrimestre 2026}
\UASetCourse{Sistemas Distribuidos}
\UASetDate{2026}


\UASetClassNumber{02}
\UASetClassTitle{RPC, timeouts, retries e idempotencia}
\UASetDocKind{Guía de ejercicios}

\title{Clase 02 --- RPC, timeouts, retries e idempotencia}
\author{\UAProfessor}
\date{\UADate}

\begin{document}

\section{Indicaciones generales}
Resuelva cada ejercicio declarando supuestos, modelo de fallas, garantía y trade-off. Cuando corresponda, construya una ejecución verificable y vincule la respuesta con evidencia observable.

\section{Nivel 1: fundamentos y reconocimiento}
\begin{uaexercise}
Definir RPC con precisión. Explicar qué intenta abstraer y por qué esa abstracción es problemática.
\end{uaexercise}

\begin{uaexercise}
Listar y explicar al menos 5 diferencias fundamentales entre una llamada local y una llamada remota.
\end{uaexercise}

\begin{uaexercise}
Describir todas las etapas de una llamada RPC e identificar en cuáles pueden ocurrir fallas.
\end{uaexercise}

\begin{uaexercise}
Construir una línea de tiempo de una llamada remota con timeout. Incluir al menos: envío del cliente, posible retransmisión o retry, vencimiento del timeout, ejecución del servidor y respuesta tardía. Indicar qué eventos pueden ocurrir después de que el cliente dejó de esperar.
\end{uaexercise}

\begin{uaexercise}
Elegir una comunicación realista y listar las capas involucradas. Para cada capa, indicar qué respuesta puede tener ante una falla y cómo esa respuesta ayuda o complica el diagnóstico.
\end{uaexercise}


\section{Nivel 2: ejecuciones y fallas}
\begin{uaexercise}
Explicar por qué la red rompe la transparencia de RPC.
\end{uaexercise}

\begin{uaexercise}
Construir un ejemplo donde una operación se ejecute pero el cliente no lo sepa.
\end{uaexercise}

\begin{uaexercise}
Definir formalmente qué es un timeout.
\end{uaexercise}

\begin{uaexercise}
Explicar por qué un timeout no implica crash.
\end{uaexercise}

\begin{uaexercise}
Diseñar un escenario donde un timeout ocurra pero la operación sea exitosa.
\end{uaexercise}


\section{Nivel 3: diseño y comparación}
\begin{uaexercise}
Diseñar un escenario donde no haya timeout pero la operación falle.
\end{uaexercise}

\begin{uaexercise}
Analizar el trade-off entre:
\begin{itemize}
\item timeouts cortos
\item timeouts largos
\end{itemize}
\end{uaexercise}

\begin{uaexercise}
Definir retry y explicar su objetivo.
\end{uaexercise}

\begin{uaexercise}
Explicar formalmente por qué retry introduce duplicación.
\end{uaexercise}

\begin{uaexercise}
Construir una línea temporal donde retry produce doble ejecución.
\end{uaexercise}


\section{Nivel 4: integración y defensa}
\begin{uaexercise}
Analizar el efecto de múltiples retries sobre el sistema.
\end{uaexercise}

\begin{uaexercise}
Comparar:
\begin{itemize}
\item sin retry
\item con retry
\end{itemize}
en términos de liveness y safety.
\end{uaexercise}

\begin{uaexercise}
Definir at-least-once, at-most-once y exactly-once.
\end{uaexercise}

\begin{uaexercise}
Dar ejemplos reales de cada semántica.
\end{uaexercise}

\begin{uaexercise}
Explicar por qué exactly-once es difícil en sistemas distribuidos.
\end{uaexercise}


\end{document}
